% ============================================================
\chapter*{Preface}
\addcontentsline{toc}{chapter}{Preface}
% ============================================================

\textbf{Operations Research} (OR) is a scientific discipline that relies on
mathematical, statistical, and algorithmic methods to support decision-making.
Born during World War~II within Allied military headquarters, it has since
become an indispensable tool in industry, logistics, finance,
telecommunications, and many other fields.

\section*{Course objectives}

This course is designed for second- and third-year undergraduate students
(L2--L3) in mathematics, computer science, or engineering.  Upon completion,
students will be able to:

\begin{itemize}
  \item Formulate real-world problems as mathematical programs (linear,
        integer, nonlinear).
  \item Solve these programs using fundamental algorithms (simplex,
        branch-and-bound, shortest path, maximum flow, etc.).
  \item Interpret results through duality and sensitivity analysis.
  \item Use software tools (Python, PuLP, SciPy, NetworkX) to solve
        realistically sized problems.
  \item Understand the principles of discrete-event simulation and queueing
        theory.
\end{itemize}

\section*{Organization of the book}

The book is organized into eleven chapters, designed to be studied
sequentially:

\begin{description}
  \item[Chapter 1 --- Introduction to OR.]
    History, definitions, and general modelling methodology.
  \item[Chapter 2 --- Linear Programming: Formulation and Geometry.]
    Building linear models, geometric interpretation, feasible regions, and
    optimal vertices.
  \item[Chapter 3 --- The Simplex Method.]
    The simplex algorithm in tableau form, initialization via the Big-$M$
    method and the two-phase method.
  \item[Chapter 4 --- LP Duality.]
    Constructing the dual, duality theorems, economic interpretation.
  \item[Chapter 5 --- Sensitivity Analysis.]
    Coefficient ranging, parametric analysis, shadow price interpretation.
  \item[Chapter 6 --- Transportation and Assignment Problems.]
    Balanced and unbalanced transportation models, stepping-stone method,
    assignment problem and the Hungarian algorithm.
  \item[Chapter 7 --- Applied Graph Theory.]
    Shortest paths (Dijkstra, Bellman-Ford), minimum spanning trees (Kruskal,
    Prim).
  \item[Chapter 8 --- Network Flows.]
    Maximum flow (Ford-Fulkerson), minimum cut, minimum-cost flow.
  \item[Chapter 9 --- Integer Programming.]
    Branch-and-bound, cutting planes, linear relaxation.
  \item[Chapter 10 --- Nonlinear Programming.]
    Karush-Kuhn-Tucker conditions, convex programming, gradient methods.
  \item[Chapter 11 --- Simulation and Queueing Theory.]
    Discrete-event simulation, M/M/1 and M/M/c models, Little's law.
\end{description}

\section*{Prerequisites}

The required background includes:
\begin{itemize}
  \item \textbf{Linear algebra}: matrices, systems of linear equations, bases,
        and rank.
  \item \textbf{Calculus}: functions of several variables, partial derivatives,
        convexity.
  \item \textbf{Probability}: standard distributions, expectation, Poisson
        processes (for Chapter~11).
  \item \textbf{Algorithms}: basic notions of complexity and Python
        programming.
\end{itemize}

\section*{Conventions and notation}

Throughout this book we use the following conventions:

\begin{center}
\renewcommand{\arraystretch}{1.3}
\begin{tabular}{cl}
  \toprule
  \textbf{Notation} & \textbf{Meaning} \\
  \midrule
  $\R$            & Set of real numbers \\
  $\R^n$          & Euclidean space of dimension $n$ \\
  $\Z$            & Set of integers \\
  $\N$            & Set of natural numbers \\
  $\mathbf{x}$    & Column vector in $\R^n$ \\
  $A \in \R^{m \times n}$ & Real $m \times n$ matrix \\
  $\mathbf{c}^\top \mathbf{x}$ & Inner (dot) product \\
  $\norm{\mathbf{x}}$ & Euclidean norm \\
  $\abs{S}$       & Cardinality of the set $S$ \\
  $\argmin$, $\argmax$ & Argument of the minimum / maximum \\
  \bottomrule
\end{tabular}
\end{center}

Theorems, definitions, and propositions are numbered by chapter.  Examples are
highlighted in boxes, and exercises appear at the end of each chapter.

\begin{remark}
\textbf{Remark} boxes provide additional insights or caveats.
\textbf{Best practice} boxes offer methodological advice for modelling and
implementation.
\end{remark}

\section*{Software used}

Numerical examples are implemented in \textbf{Python~3} using the following
libraries:

\begin{itemize}
  \item \texttt{scipy.optimize.linprog} --- solving linear programs;
  \item \texttt{PuLP} --- modelling and solving LP/ILP;
  \item \texttt{NetworkX} --- graphs, shortest paths, flows;
  \item \texttt{SimPy} --- discrete-event simulation;
  \item \texttt{NumPy}, \texttt{Matplotlib} --- numerical computing and
        visualization.
\end{itemize}

\begin{codeblock}[title=\textbf{Installing dependencies}]
\begin{minted}{python}
# pip install numpy scipy matplotlib pulp networkx simpy
import numpy as np
from scipy.optimize import linprog
import pulp
import networkx as nx
\end{minted}
\end{codeblock}

\section*{How to use this course}

Each chapter follows a recurring structure:
\begin{enumerate}
  \item \textbf{Motivation} --- a concrete problem that introduces the
        theoretical need.
  \item \textbf{Theory} --- definitions, theorems, and proofs.
  \item \textbf{Algorithms} --- formal description and worked example.
  \item \textbf{Implementation} --- commented Python code.
  \item \textbf{Exercises} --- problems of increasing difficulty
        ($\star$ to $\star\star\star$).
\end{enumerate}

\begin{bestpractice}[title=\textbf{Active learning}]
Operations research is best learned by \textit{doing}. We strongly recommend:
\begin{itemize}
  \item Working through examples by hand before consulting the solution.
  \item Implementing algorithms in Python.
  \item Attempting exercises before looking at solutions.
  \item Exploring variants of the proposed problems.
\end{itemize}
\end{bestpractice}

\section*{Brief history}

\begin{itemize}
  \item \textbf{1937}: Kantorovich formulates the first linear programming
        problems for economic planning.
  \item \textbf{1939--1945}: British and American OR teams optimize military
        operations (convoys, radar, bombing campaigns).
  \item \textbf{1947}: George Dantzig publishes the simplex method.
  \item \textbf{1951}: Kuhn and Tucker state optimality conditions for
        nonlinear programming.
  \item \textbf{1956}: Ford and Fulkerson propose the maximum-flow algorithm.
  \item \textbf{1958}: Gomory introduces cutting planes for integer
        programming.
  \item \textbf{1960}: Land and Doig formalize branch-and-bound.
  \item \textbf{1979}: Khachiyan shows LP is solvable in polynomial time
        (ellipsoid method).
  \item \textbf{1984}: Karmarkar publishes the interior-point method.
  \item \textbf{1990s--2020s}: Development of commercial solvers (CPLEX,
        Gurobi) and open-source solvers (GLPK, CBC); widespread adoption in
        industry.
\end{itemize}

\section*{Main references}

\begin{enumerate}
  \item \textsc{Hillier} F.S., \textsc{Lieberman} G.J.,
        \textit{Introduction to Operations Research}, McGraw-Hill,
        11th~ed., 2021.
  \item \textsc{Bertsimas} D., \textsc{Tsitsiklis} J.N.,
        \textit{Introduction to Linear Optimization}, Athena Scientific, 1997.
  \item \textsc{Wolsey} L.A.,
        \textit{Integer Programming}, Wiley, 2nd~ed., 2020.
  \item \textsc{Chv\'atal} V.,
        \textit{Linear Programming}, W.H.~Freeman, 1983.
  \item \textsc{Winston} W.L.,
        \textit{Operations Research: Applications and Algorithms},
        Cengage, 4th~ed., 2003.
  \item \textsc{Bazaraa} M.S., \textsc{Jarvis} J.J., \textsc{Sherali} H.D.,
        \textit{Linear Programming and Network Flows}, Wiley, 4th~ed., 2010.
\end{enumerate}

\section*{Acknowledgements}

I thank the colleagues and students who contributed to improving this course
through their comments and suggestions.  Any remaining errors are my sole
responsibility.

\vfill
\begin{flushright}
\textit{Happy reading and happy optimizing!}
\end{flushright}
